% =====================================================================
\section{Demo, hạn chế và kết luận}
\label{sec:conclusion}

\subsection{Demo}
\label{sec:demo}

\file{app/demo_app.py} là một giao diện Gradio, chạy được trên GPU 8\,GB. Người dùng nhập truy vấn tiếng Anh, chọn CLIP thuần
hoặc CLIP + đồ thị, rồi xem top-$k$ ảnh. Mỗi ảnh kèm bảng giải thích như Mục~\ref{sec:method-expl}, khung của các vật thể khớp
vẽ trên ảnh, và cột ``bỏ phần này'': chấm lại ảnh bằng chính mô hình sau khi bỏ phần khớp đó khỏi đồ thị ảnh. Câu gõ vào được
tách bằng FACTUAL-T5 và cùng luật số lượng như lúc đóng gói dữ liệu. Phép kiểm tra ở Mục~\ref{sec:checks} bảo đảm demo cho
cùng kết quả với pipeline đánh giá. Video dự phòng theo kịch bản \file{docs/demo-video-script.md}.

\subsection{Hạn chế}

\begin{itemize}
\item \textbf{Chỉ 3 seed.} Mọi kiểm định giữa các cấu hình chỉ dựa trên 3 seed mỗi bên.
\item \textbf{Val được dùng nhiều lần} để chọn từng bước của mô hình chính VG, nên số val hơi lạc quan (từ val sang test giảm
0,96 điểm R@1). Test vẫn chỉ chấm một lần.
\item \textbf{Chưa tách được đóng góp của từng loại cạnh} trong mô hình chính VG: bước đồ thị liền gộp bốn thay đổi, và chưa có
ablation chỉ bỏ cạnh kNN.
\item \textbf{Đếm số lượng coi mọi người là một loại.} Câu ``a lone woman'' trong ảnh nhiều người bị coi là sai số lượng, và
bảng $g$ đôi khi kéo phần khớp sang một node khác nhãn.
\item \textbf{Lỗi parser.} Thư viện FactualSceneGraph 0.7.3 cắt token cuối của câu dài nhất trong mỗi batch. Khoảng 1,8\%
caption bị ảnh hưởng; 0,26\% (train), 0,24\% (val) và 0,31\% (test) mất một phần chữ. Lỗi xảy ra như nhau ở cả ba split.
Nhóm không parse lại dữ liệu, giữ nguyên các số đã báo, và đã sửa lỗi này trong demo.
\item \textbf{Thuộc tính và chữ trong ảnh chưa được dùng}, dù nhiều ca thất bại liên quan tới chúng.
\item \textbf{Giới hạn của xếp lại.} 3,3\% câu test VG có ảnh đúng nằm ngoài top-50 của CLIP và không cứu được bằng xếp lại.
\item \textbf{MKG-W.} Phương pháp năm 2024 không đóng góp vào dòng chính ($\beta = 0$). Đặc trưng ảnh của NativE (383 chiều,
do tác giả trích sẵn) khác CLIP ViT-B/32 của nhóm. Cache Wikidata phải lấy lại ngày 17/11: 441 thực thể mất nhãn tiếng Anh,
số truy vấn val khớp được thực thể đổi từ 704 xuống 703. Bộ câu tự viết chỉ có 25 câu.
\end{itemize}
Hạn chế riêng của phần đánh giá lời giải thích và phần MMKG tiếng Việt ghi ở Mục~\ref{sec:xai} và Mục~\ref{sec:vnmmkg}.

\subsection{Kết luận}

Báo cáo trả lời ba câu hỏi của đề như sau.
\begin{enumerate}
\item \textbf{Kết quả nào phù hợp nhất?} Trên VG $\cap$ COCO, xếp lại top-50 của CLIP bằng scene graph nâng R@1 test từ
39,32 lên 47,95. Trên MKG-W, đếm cạnh tri thức nâng R@1 test từ 24,89 lên 29,82.
\item \textbf{Tri thức quan hệ giúp ở đâu?} Ở câu có nhiều vật thể và có quan hệ: phần lời tăng theo số vật thể, còn câu chỉ
có một vật thể thì không được lợi. Ở câu gián tiếp: trên 15 câu tự viết, R@1 tăng từ 40 lên 80, trong khi câu trực tiếp
không đổi. Các phép đối chứng cho thấy phần lời đến từ đúng cấu trúc: thay đồ thị của ảnh khác hoặc xáo cạnh tri thức đều
làm kết quả tụt dưới CLIP.
\item \textbf{Vì sao mô hình trả về kết quả này?} Lời giải thích là chính các đại lượng của điểm đồ thị. Với mô hình chính
của VG, điểm tách chính xác thành từng phần. Mục~\ref{sec:xai} đo mức trung thực của lời giải thích đó.
\end{enumerate}
Nhóm cũng ghi lại các kết quả âm có bằng chứng: ở cấp~1 cạnh đóng góp ít, và NativE không thêm được gì khi đã có kênh đếm
cạnh.
